% !Mode:: "TeX:UTF-8"
% 哈尔滨校区本科生开题报告正文
% 题目：基于记忆与想像建模的视觉-语言目标跟踪方法研究

\section{课题来源及研究的目的和意义}

\subsection{课题来源}

本课题来源于国家自然科学基金青年基金项目“基于自然语言查询的视觉目标跟踪方法研究”（项目负责人：马丁）。该项目面向自然语言查询条件下的视觉目标跟踪问题开展应用研究；本课题作为其支撑课题，针对视觉-语言跟踪中时序建模薄弱的关键问题，研究基于记忆与想像建模的视觉-语言目标跟踪方法。

\subsection{课题背景及研究目的和意义}

视觉目标跟踪是计算机视觉领域的基础任务，在智能监控、自动驾驶、人机交互等场景中有着广泛应用\cite{fan2021lasot}。近年来，随着深度学习与跨模态预训练技术的成熟，跟踪任务的输入形式从传统的“首帧给定目标框”扩展为“首帧目标框+自然语言描述”，甚至仅凭一段自然语言描述即可在视频中锁定目标，由此形成了视觉-语言目标跟踪这一新兴研究方向\cite{wang2021tnl2k,guo2024divert}，其输入输出形式如图~\ref{fig:vlttask}~所示。语言描述的引入使模型能够理解“跟踪谁、跟踪目标的哪一部分”，显著提升了人机协同的灵活性，但同时也对模型的多模态理解与时序建模能力提出了更高要求。

% 图 1-1：视觉-语言目标跟踪任务示意
\begin{figure}[htbp]
  \centering
  \includegraphics[width=0.75\textwidth]{fig1-1-vlt-task}
  \caption{视觉-语言目标跟踪任务示意}
  \label{fig:vlttask}
\end{figure}

然而，现有视觉-语言目标跟踪方法普遍沿用“逐帧跨模态匹配”的建模范式，时序建模能力薄弱、过度依赖单帧瞬时特征，难以充分利用视频时序关联信息：模型在每个时刻独立融合当前帧观测与语言描述，帧间仅靠模板替换或短窗口特征拼接维持弱关联，因而在目标遮挡、形态变化及复杂背景干扰场景下易出现跟踪漂移、目标丢失等现象，无法满足长时序、高稳定性的跟踪需求\cite{hu2023mgit,ma2023decouple}。与此同时，LaSOT、TNL2K、MGIT等评测基准的序列长度已从数百帧扩展到数千帧乃至超长视频，评测形态正向长时序、多轮交互演进；MGIT的实验进一步表明，多数多模态跟踪器在给定完整长文本标注时性能反而下降，瓶颈同时出现在文本侧的长上下文处理能力与视觉侧的时序建模能力\cite{fan2021lasot,wang2021tnl2k,hu2023mgit}。已有工作分别从模板动态更新与状态空间压缩两个方向逼近该问题：前者更新的只是语义描述，本质上仍属“用当前观测覆盖过去”；后者把历史压缩进隐式演化的状态空间，虽已联合建模视觉与语言参考，却缺乏可检索、可校验的显式记忆结构，也未对目标未来状态作出显式预判，因而长时序跟踪漂移与丢失问题尚未得到真正解决\cite{zhao2025selfadaptive,ma2023decouple,liu2025mambavlt}。

为此，本课题借鉴人类视觉认知的记忆与推演机制——工作记忆缓存短时上下文，情景记忆留存历史目标信息，内部认知模型对目标后续状态作出预判——构建融合记忆留存与未来状态想像的全时序建模跟踪框架。具体而言：建立视觉-语言时序特征编码方式，为时序建模提供表达充分、时序可分辨的特征基础；构建感知—认知记忆库及其自适应更新机制，实现历史有效特征的筛选、存储与复用；在潜空间中对目标未来状态进行去噪式推演，对目标位置与外观演化作出合理预判，使跟踪器由“被动响应”转向“主动预判”；最终将当前帧特征、记忆特征与推演特征有效融合，输出时序一致性更强的跟踪预测。通过挖掘历史有效特征、预判目标动态变化，本课题有望提升模型在复杂场景下的跟踪精度与鲁棒性，改善长时序跟踪中的漂移与丢失问题。

在理论层面，本研究针对传统跟踪方法重空间、轻时序的建模局限，将跟踪过程由逐帧的判别问题转化为跨时间的状态演化问题，补充了多模态时序感知体系，探索了语言语义先验与视觉时序信息在长时序条件下的协同规律，为视觉-语言跟踪任务提供了新的建模思路与理论参考。

在应用层面，该方法适配智能监控、机器人视觉、人机交互等主流智能感知场景，能够有效解决传统算法稳定性差、泛化能力弱的工程痛点，为智能视觉设备的落地应用与性能优化提供可靠技术支撑。综上，本课题面向视觉-语言目标跟踪中时序建模薄弱的现实问题，以记忆与想像建模为核心手段，兼具理论探索价值与工程应用前景，具备重要的研究价值与广阔的应用前景。

\section{国内外在该方向的研究现状及分析}

\subsection{视觉-语言目标跟踪研究现状}

视觉-语言目标跟踪源于纯视觉跟踪在参考信息上的固有局限：早期方法以首帧目标框为唯一参考，缺乏高层语义约束，外观剧变或相似物干扰时易产生歧义，相关综述对该方向的技术演进脉络有系统梳理\cite{ma2024review}。Li等人\cite{li2017tracking}最早以自然语言引导跟踪并构建OTB99\_Lang评测集，验证了语言先验的有效性；Feng等人\cite{feng2021siamnlt}提出孪生自然语言跟踪器，首次将语言模态直接接入孪生跟踪框架；Wang等人\cite{wang2021tnl2k}构建TNL2K基准并引入对抗样本与模态切换挑战；Guo等人\cite{guo2024divert}以纯卷积模态混合器将LaSOT成功率由50.7\%提升至65.2\%，证明语言表征的价值不依赖Transformer架构。此后，评测基准沿“规模—难度—语义粒度”持续扩张，MGIT\cite{hu2023mgit}以150段超长视频、约203万帧按动作、活动、故事三种语义粒度标注，长时跟踪对时序建模提出了远高于短时场景的要求。

在融合方式上，早期工作普遍采用“先定位、再跟踪”的两阶段结构：JointNLT\cite{zhou2023jointnlt}将视觉定位与跟踪重构为统一任务并引入时序建模模块，All in One\cite{zhang2023allinone}以统一Transformer骨干实现特征提取与跨模态交互一体化。此类方法跨模态对齐显式、可解释性强，但定位误差会直接传播至跟踪环节，历史信息仅作为辅助线索。另一条路线关注“如何让语言描述跟上目标变化”：CiteTracker\cite{li2023citetracker}将目标图像块转为描述性文本并在线更新，DecoupleTNL\cite{ma2023decouple}以长短期调制整合文本与片段上下文，ATCTrack\cite{feng2025atctrack}将目标线索与动态目标状态对齐；但它们的作用对象始终是“参考信息”本身，未形成对历史状态独立的留存与检索机制。

近期，长序列建模成为新的增长点：MambaVLT\cite{liu2025mambavlt}以时变状态空间模型线性建模长序列上下文，并以模态选择模块动态权衡视觉与语言参考；ODTrack\cite{zheng2024odtrack}通过在线令牌传播维持帧间关联；但二者传递的都是隐式演化的历史状态，缺乏可检索、可遗忘的显式记忆结构，目标未来状态亦未进入建模范围。综上，现有视觉-语言跟踪方法在跨模态对齐与语义演化上已相对成熟，但时序维度仍存在三层局限：观测的瞬时性，遮挡或形变时缺乏可靠历史参照；记忆的缺失，缺少可检索、可遗忘、可校验的显式记忆结构；预测的被动性，未显式建模目标未来状态，遮挡期只能被动等待全局重检测。

\subsection{记忆机制研究现状}

记忆机制研究横跨认知心理学与神经计算两个源头：Atkinson-Shiffrin模型将人类记忆划分为感觉、短时与长时记忆并指出“巩固”过程\cite{atkinson1968human}，神经图灵机与可微神经计算机提供了“可微、可读写、可寻址”的外部记忆机制\cite{graves2014ntm,graves2016hybrid}。视频目标分割是其最成功的应用场域：STM\cite{oh2019stm}将历史帧特征以键值形式存入记忆库、以注意力读出，奠定“记忆写入—记忆读出”范式；XMem\cite{cheng2022xmem}构建感觉、工作与长时记忆三种时间尺度的存储结构，并设计巩固、增强与遗忘机制。其对本课题的启示在于：分层记忆优于单一记忆，“巩固—增强—遗忘”是记忆系统的必要三要素；而其记忆内容全为单模态视觉特征，尚未在存储层引入语言语义作为写入条件、检索键与遗忘判据，这是一条尚未被占据的差异化路径。

记忆机制在视觉跟踪中的引入更晚也更浅：HIPTrack\cite{cai2024hiptrack}以历史提示缓解单模板的表示瓶颈，另有工作对历史帧评分择优写入以缓解误差累积\cite{zheng2025robust}，其背后是“记忆污染”问题——错误特征写入后在后续帧中被反复强化、难以回溯修正。综合来看，现有记忆机制在视觉跟踪场景下仍存在三个未解问题：内容单模态，记忆条目未与语言语义绑定组织，语言不参与记忆的检索与遗忘；更新粗放，跟踪中的记忆写入多为固定频率更新或单一评分择优，缺少面向语言条件的多因子自适应更新与显式遗忘机制；与预测解耦，记忆既未用于未来状态推演，也缺乏反向校验记忆质量的闭环。

\subsection{目标未来状态预测研究现状}

目标未来状态的预测性建模主要源于世界模型研究：Ha与Schmidhuber\cite{ha2018worldmodels}首次证明无需显式建模物理规律即可从观测中习得环境动力学，确立了世界模型范式；Dreamer\cite{hafner2020dreamer}证明可在潜空间中“想像”未来轨迹并反向传播学习，后续工作以统一超参数验证了该范式的稳定性\cite{hafner2025mastering}。其确立的潜空间动力学、想像展开与预测误差最小化等概念，构成本课题“想像推演”模块的方法论基础；但此类工作以动作为条件，想像服务于动作选择而非目标状态估计。

在实现空间上，趋势已由像素空间转向表征空间：I-JEPA\cite{assran2023ijepa}主张在抽象表征空间预测以避开不可预测细节，相关工作以冻结的视觉基础模型编码器加轻量级动力学预测器在潜空间自回归预测未来特征；去噪扩散模型及其潜空间变体确立了“去噪即生成”路径\cite{ho2020ddpm,rombach2022ldm}，为“以当前观测为条件、从噪声中恢复未来状态”提供机制支撑。

视觉跟踪场景中，未来状态预测刚刚起步：CST-WM\cite{hu2026cstwm}面向具身跟踪提出因果结构世界模型，揭示动作条件下的“因果幻觉”现象，主张预测结构必须与任务结构一致；LeWorldModel\cite{maes2026leworldmodel}证明联合嵌入预测架构可在紧凑潜空间中稳定训练，为轻量级推演器提供实现模板；DeltaWorld\cite{kerssies2026delta}以差分令牌压缩相邻帧变化，表明“只预测变化”更契合在线跟踪约束；已有工作将目标指定扩展为文本、点、框与掩码的统一提示，并以潜世界模型预测未来表征提升时序鲁棒性\cite{xie2026uss}。但上述工作的推演仍以动作或空间提示为条件，据本文调研范围所限，语言指代条件下的目标状态推演尚未被系统研究。

值得关注的是，记忆与预测长期被作为两个独立问题处理，而视觉-语言-动作模型领域近期出现了联合建模的工作：MemoryVLA++\cite{shi2026memoryvla}以双流感知-认知记忆库加视频扩散世界模型，在潜空间推演未来表征并与记忆增强表征融合，验证了记忆与想像联合建模在长时程任务上的增益，是本课题最直接的方法学参照；但它面向机器人操作、以动作与指令为条件，且依赖7B视觉-语言模型与1.5B参数视频扩散世界模型生成像素级未来观测，模型规模与显存开销远超在线跟踪的逐帧预算。综合而言，世界模型迁移至视觉-语言跟踪面临四重障碍：条件化不匹配，需以语言与记忆替代动作为条件；幻觉与漂移风险，一次幻觉即可能跟丢；实时性冲突，多步去噪与显存开销难以满足在线逐帧预算；监督信号缺失，仅有目标框标注。

\subsection{研究现状分析与存在的问题}

综合三条技术主线：跨模态对齐与语义演化、记忆管理系统、潜空间预测均已分别形成较成熟的方法体系，但尚未在视觉-语言跟踪任务上交汇——视觉-语言跟踪的时序建模停留在模板更新层面，记忆机制的内容为单模态视觉特征且未进入未来推演环节，世界模型以动作为条件、面向具身决策。据本文调研范围所限，面向语言指代目标的状态推演仍属空白。

据此，本课题提炼出三个待解决的关键问题：第一，如何构建语言语义参与的多模态记忆，使记忆的写入、检索与遗忘均由视觉-语言联合信息驱动；第二，如何以记忆与语言为条件在潜空间中推演目标未来状态并抑制幻觉，建立预测-校正机制保证推演可靠；第三，如何把推演得到的未来表征与当前帧表征有效融合并送入预测环节，使跟踪同时利用“现在”与“将要”。上述三个问题分别对应本课题的记忆库构建与自适应更新、记忆条件下的世界模型潜空间想像、以及未来表征与当前表征的融合注入与跟踪预测三项研究内容。

基于上述分析，本课题选择无模板视觉-语言跟踪方法MVLM\cite{park2026mvlm}作为基线。MVLM\cite{park2026mvlm}仅以自然语言作为语义锚点，不使用首帧给定的目标模板，并已通过temporal memory参与置信度估计与局部搜索/全局重定位门控。因此，本文不将其视为“没有记忆”的方法，而是进一步关注其时间记忆在目标内容组织、显式检索和未来状态预测方面的扩展空间。该基线具有清晰的置信度建模和搜索区域调度机制，便于在统一实验口径下分析内容记忆与前瞻推演的增量作用。在方法上，本课题借鉴MemoryVLA++\cite{shi2026memoryvla}中记忆库与世界模型的设计，并将其由动作条件、像素生成的设定，改造为面向在线跟踪的“语言与记忆双条件、潜空间表征级推演”的设定。具体研究内容见第\ref{sec:content}~章。

\section{主要研究内容}
\label{sec:content}

针对上一章提炼的三个关键问题，本章以无模板视觉-语言跟踪方法MVLM\cite{park2026mvlm}为基线，从记忆与想像两个方向对其增强，给出三项研究内容、总体研究框架与研究内容分解。

\subsection{基线模型MVLM及其记忆机制边界}

本课题以无模板（template-free）视觉-语言跟踪方法MVLM\cite{park2026mvlm}作为基线。该方法采用冻结的CLIP文本编码器提取语言表征，并按照官方配置使用视觉骨干网络提取图像表征，将视觉token与语言token拼接为统一序列送入Transformer联合编码；再由分类、尺寸回归与偏移回归三个分支预测头在特征网格上逐点生成候选框集合，经Top-K选择与BBOX过滤得到最终预测。基线以区域平均的视觉-语言相关度作为定位依据，并以相关度边际（CM）损失拉大真值区域与最难背景之间的相关度差。

% 图 3-1：基线 MVLM 跟踪框架
\begin{figure}[htbp]
  \centering
  \includegraphics[width=0.95\textwidth]{fig3-1-mvlm}
  \caption{基线MVLM跟踪框架}
  \label{fig:mvlm}
\end{figure}

在置信度层面，基线由相关度与分类得分的边际构成瞬时置信度，再经指数滑动平均得到时序平滑的置信度，据其在“局部ROI搜索”与“全局重定位”两种策略间在线切换；理论上，基线将跟踪失败分解为ROI排除与ROI内误定位两项，并给出概率上界。基线的整体架构与置信度驱动的跟踪流程如图~\ref{fig:mvlm}~所示。

需要区分的是，MVLM已经引入temporal memory，但其记忆对象是标量形式的置信度序列，经指数滑动平均后仅用于判断“当前跟踪是否可靠”，并不保存任何目标外观与语义内容；且其语言信息以全局单一语义向量的形式参与相关度计算，未与具体记忆条目建立绑定关系。因此，其现有记忆主要回答“当前跟踪是否可靠”，而不是完整保存和复用目标的历史内容。基于这一边界，本课题拟解决三方面问题：其一，现有时间记忆尚未形成视觉-语言联合的显式内容记忆，难以对目标外观、语义和时序状态进行独立存储、检索与遗忘；其二，现有机制主要用于当前可靠性判断，尚未显式预测目标未来状态；其三，历史记忆与未来预期尚未直接进入候选框的生成和筛选过程。

\subsection{研究内容一：面向视觉-语言跟踪的内容记忆库构建与自适应更新}
\label{sec:mem}

针对上述内容记忆能力不足的问题，本课题在MVLM已有时间记忆机制的基础上，进一步构建键值对（key-value）形式的视觉-语言内容记忆库。记忆条目采用感知—认知双流结构：感知流以外观表征为键，认知流以语义表征为键，二者的值统一存储对应的时序状态信息；条目附带时间位置编码，以注意力机制检索读出，容量受限时按冗余度淘汰旧条目，选型思路参考XMem\cite{cheng2022xmem}、STM\cite{oh2019stm}等视频记忆机制。

在框架中，记忆读出携带“历史目标的外观、语义与状态信息”，作为世界模型潜空间推演的条件输入；记忆写入以跟踪置信度与推演一致性双重门控，仅在当前帧证据可靠时允许入库，以抑制错误特征写入；语言描述变更或跟踪目标切换时重置记忆，以适配MVLM“改语言即换目标”的无模板特性。由此，已有时间记忆负责跟踪可靠性评估，新增内容记忆负责目标历史信息的显式保存与复用，二者在功能上相互补充。

本部分预期产出内容记忆库的结构设计、有无内容记忆的对照实验方案，以及记忆质量的度量方式。

\subsection{研究内容二：记忆条件下的世界模型潜空间想像}

本部分为核心增量，针对决策无前瞻与候选框生成封闭于当前帧的不足，负责产出可用于注入的双重前瞻证据：未来帧的目标表征，以及每个候选框的事前可靠度预判。世界模型采用轻量Transformer结构，在低维潜空间做表征级推演以满足在线跟踪的实时约束；推演条件为三路信息：记忆库读出、语言语义向量与当前帧视觉特征。

预测头A承担视觉前瞻，输出“交互/运动后的视觉”，即未来帧的目标表征。该表征指向研究内容三的注入环节——在遮挡、快速运动等当前帧证据缺失的情形下，由假想观测补全证据，引导候选框在正确位置产生。

预测头B承担权值前瞻，为每个候选框预测权值，即对“画出来的框有多可靠”的事前预判。该权值指向研究内容三的筛选环节，用于Top-K选择与BBOX过滤，权值高的候选才进入最终预测，使候选框筛选由“只看当前帧”变为“当前证据与未来预期”联合决定。

本部分预期产出推演器的结构与训练目标、双重前瞻预测方案与幻觉度量指标。上述训练目标的可行性依据在于：现有数据集的时序边界框标注已足以构造监督信号，不需要新增人工标注；且以单步推演替代多步自回归展开，可在机制上控制误差累积。具体的网络结构、损失权重与推理门控实现见第~\ref{sec:worldmodel}~节。

\subsection{研究内容三：未来表征与当前表征的融合注入与跟踪预测}

本部分是新增模块与基线的直接接口。参照系对齐方面，未来表征描述未来时刻的目标状态，须与当前帧空间网格对齐后才能参与候选框生成；注入方式比较拼接与交叉注意力两条路线，并界定预测头的参数适配与微调范围，视觉主干全程冻结。

权值引导的筛选方面，以前瞻权值与基线分类得分、相关度得分联合排序，确定加权Top-K选择、BBOX过滤与最终预测流程，使“生成哪些框、保留哪些框”由当前证据与未来预期共同决定。

需要说明的是，基线的置信度聚合与“局部搜索、全局重定位”门控保持原有形式不变，主要改动集中于候选框环节；在保持数据划分、评价协议、主干网络和搜索策略一致的条件下，通过逐模块消融实验分析新增内容记忆与前瞻机制的独立贡献。

\subsection{总体研究框架与研究内容分解}
\label{sec:decomp}

本课题以“记忆留存历史、想像推演未来”为主线，构建“观测—记忆—想像—注入”的全时序建模框架：搜索图像与语言描述经统一编码后，由内容记忆库留存目标历史，由世界模型在潜空间推演未来；双重前瞻证据（未来表征与候选框权值）注入候选框的生成与筛选，最终预测仍由基线三分支头完成；置信度聚合与搜索区域门控保持基线形式，负责搜索区域的在线调度。

加入记忆库与世界模型后的整体跟踪框架如图~\ref{fig:framework}~所示，其中记忆库负责留存目标历史，世界模型负责在潜空间推演未来状态，二者的输出经注入与筛选环节共同作用于候选框生成。
\begin{figure}[htbp]
  \centering
  \includegraphics[width=0.95\textwidth]{fig3-2-framework}
  \caption{加入记忆库与世界模型的总体跟踪框架}
  \label{fig:framework}
\end{figure}

三项研究内容与三个关键问题一一对应：内容记忆库回应“记不住”，世界模型回应“想不到”，融合注入与权值筛选回应“用不上”，共同支撑“基于记忆与想像建模”的核心思路，具体对应关系如表~\ref{tab:t3-1}~所示。

\begin{table}[htbp]
  \caption{基线结构性不足与研究内容对应关系}
  \label{tab:t3-1}
  \centering
  \begin{tabularx}{\textwidth}[t]{A{0.32} A{0.68}}
    \toprule
    结构性不足 & 对应研究内容 \\
    \midrule
    记忆无内容 & 研究内容一：内容记忆库构建与自适应更新 \\
    决策无前瞻 & 研究内容二：记忆条件下的世界模型潜空间想像 \\
    候选框生成封闭于当前帧 & 研究内容三：未来表征与当前表征的融合注入与跟踪预测 \\
    \bottomrule
  \end{tabularx}
  \tabdepthfix
\end{table}

\section{研究方案}

本章给出本课题的研究方案与实验设计。本课题前期已完成相关文献的系统梳理与储备（参考文献37篇），并对基线方法MVLM的技术机制、记忆机制边界及其扩展空间进行了深入分析，明确了其时间记忆仅作用于决策层、未形成显式内容存储与检索结构的局限，为本课题的方案设计提供了直接依据；同时已完成开题报告撰写环境与图表素材的整理工作。

\subsection{总体技术路线}
\label{sec:route}

研究按三个阶段推进：阶段零复现基线并搭建评测脚本；阶段一冻结基线主干，落地内容记忆库与世界模型；阶段二解冻预测头适配层联合微调并开展消融实验。每个阶段设置明确的验收标准：阶段零以复现对齐为准，阶段一、二以模块逐个上线的增量指标为准，保证问题可定位、进度可检验。总体技术路线与阶段推进安排如图~\ref{fig:roadmap}~所示。

% 图 4-1：总体技术路线与阶段推进
\begin{figure}[htbp]
  \centering
  \includegraphics[width=0.98\textwidth]{fig4-1-roadmap}
  \caption{总体技术路线与阶段推进}
  \label{fig:roadmap}
\end{figure}

\subsection{基线复现与评测搭建}
\label{sec:repro}

按原文超参数与实现细节复现MVLM，包括冻结的CLIP文本编码器与视觉骨干网络、三分支预测头，具体骨干结构与预训练权重以官方配置为准；评测覆盖TNL2K、LaSOT、OTB99\_Lang与MGIT四个基准，采用PRE、AUC与NPR指标。复现结果与论文报告值的差距控制在一个百分点以内方可进入阶段一。

需要说明的是，基线的置信度聚合与“局部搜索、全局重定位”门控在全部实验中保持基线实现不动，作为贯穿始终的公平对照。

\subsection{内容记忆库的实现}

按第~\ref{sec:mem}~节的选型实现键值对内容记忆库：感知—认知双流条目（感知流以外观表征为键、认知流以语义表征为键，值统一存时序状态）附带时间位置编码，以注意力机制检索读出，容量受限时按冗余度淘汰旧条目。记忆特征取自视觉编码器输出层，读出结果作为世界模型的推演条件；语言描述变更时重置记忆，实现“改语言即换目标”的跟踪。

\subsection{世界模型潜空间推演器的实现}
\label{sec:worldmodel}

推演器采用轻量Transformer结构，在低维潜空间做表征级推演，层数、隐藏维与注意力头数等实现参数以满足在线实时约束为准。预测头A在潜空间预测未来帧目标表征，预测头B为每个候选框预测权值。

训练目标利用现有数据集的时序边界框标注构造，不需要新增人工标注：视觉前瞻以未来帧真值框内特征为回归目标；权值前瞻以下一帧候选框的实际重叠度与得分为回归目标。推理时以推演结果与实际观测的余弦相似度为门控，一致性低时对前瞻证据降权或弃用，抑制幻觉。

\subsection{候选框生成与筛选的改造}

未来表征经参照系对齐后注入跟踪头，注入方式比较拼接与交叉注意力两条路线；权值引导筛选以前瞻权值与基线分类得分、相关度得分联合排序，实现加权Top-K选择与BBOX过滤。视觉主干全程冻结，仅训练新增模块与预测头适配层。

\subsection{数据集、评价指标与实验环境}

训练使用TNL2K、LaSOT与VastTrack\cite{peng2024vasttrack}，与基线保持同一口径；评测使用TNL2K、LaSOT、OTB99\_Lang与MGIT，指标为PRE、AUC与NPR。同时按遮挡、快速运动、背景干扰与尺度变化构建属性子集，用于属性段分析。实验环境的硬件配置、深度学习框架与训练超参数将参照MVLM原文设置，在基线复现阶段（第~\ref{sec:repro}~节）确定；新增模块的训练安排见第~\ref{sec:stagetrain}~节分阶段训练策略。

\subsection{分阶段训练策略}
\label{sec:stagetrain}

阶段一冻结基线主干，仅训练内容记忆库、世界模型与预测头适配层，训练目标为基线跟踪损失与推演回归损失的加权和，损失权重通过验证集确定；阶段二解冻预测头适配层联合微调，并对注入方式与筛选策略做出最终选择。

\subsection{实验设计与消融方案}
\label{sec:ablation}

主实验在四个评测基准上与两类方法比较：一类为同时使用边界框与语言的跟踪方法（JointNLT\cite{zhou2023jointnlt}、UVLTrack\cite{ma2024uvltrack}等），一类为纯语言跟踪方法（QueryNLT\cite{shao2024querynlt}、MambaVLT\cite{liu2025mambavlt}等）。

组件消融遵循单变量原则，全部自同一阶段一模型出发，分四组展开：记忆库组比较无记忆、有记忆与不同淘汰策略；世界模型组比较无推演、仅视觉前瞻、仅权值前瞻与双头联合；注入方式组比较拼接与交叉注意力；筛选策略组比较基线原始Top-K与权值引导筛选，各项配置如表~\ref{tab:t4-1}~所示。

\begin{table}[htbp]
  \caption{组件消融实验矩阵}
  \label{tab:t4-1}
  \centering
  \begin{tabularx}{\textwidth}[t]{A{0.19} A{0.52} A{0.29}}
    \toprule
    消融组 & 配置 & 验证目的 \\
    \midrule
    M 记忆库 & M0 无 / M1 有 / M2 不同淘汰策略 & 内容记忆的贡献 \\
    W 世界模型 & W0 无 / W1 仅视觉前瞻 & 双头前瞻的贡献 \\
               & W2 仅权值前瞻 / W3 双头联合 & \\
    I 注入方式 & I1 拼接 / I2 交叉注意力 & 注入方式的选择 \\
    F 筛选策略 & F0 基线原始 Top-K / F1 权值引导 & 权值前瞻的落地价值 \\
    \bottomrule
  \end{tabularx}
  \tabdepthfix
\end{table}

在消融之外开展两类分析：属性子集分析按遮挡、快速运动、背景干扰与尺度变化分别统计增益分布，预期前瞻证据在遮挡与快速运动段增益最大；效率分析统计参数量与推理速度并与基线比较，验证轻量化改造满足在线实时约束。上述设计与第~\ref{sec:content}~章提炼的三个关键问题一一闭环：记忆库组回应“记不住”，世界模型组回应“想不到”，筛选策略组回应“用不上”。

\section{进度安排，预期达到的目标}

本课题研究周期为2026年9月至2027年5月，按第~\ref{sec:route}~节的三个阶段推进，进度安排如表~\ref{tab:t5-1}~所示。

\begin{table}[htbp]
  \caption{研究进度安排}
  \label{tab:t5-1}
  \centering
  \small
  \begin{tabularx}{\textwidth}[t]{A{0.14} A{0.22} A{0.38} A{0.26}}
    \toprule
    阶段 & 时间 & 主要工作 & 阶段成果 \\
    \midrule
    阶段零：基线复现 & 2026.09--2026.10 & 文献调研与实验环境搭建；复现MVLM基线，搭建评测与统计脚本 & 基线复现指标与原文对齐，形成可复用评测流程 \\
    阶段一：模块落地 & 2026.10--2027.01 & 实现内容记忆库与世界模型潜空间推演器；冻结主干训练，完成未来表征注入与权值引导筛选 & 含记忆与前瞻模块的跟踪模型，主实验初步结果 \\
    阶段二：消融优化 & 2027.01--2027.03 & 解冻预测头适配层联合微调；开展组件消融、属性子集与效率分析，确定最终配置 & 完整消融结果与四个基准的最终指标 \\
    结题撰写 & 2027.03--2027.05 & 整理实验结果，撰写毕业论文，修改定稿并准备答辩 & 毕业论文定稿与答辩材料 \\
    \bottomrule
  \end{tabularx}
  \tabdepthfix
\end{table}

预期达到的目标包括以下三个方面。

功能目标：预期实现一个记忆与想像联合增强的无模板视觉-语言跟踪方法，使其具备三项能力：其一，语言语义条件驱动的目标历史内容记忆与检索能力；其二，记忆条件下的目标未来状态潜空间推演能力；其三，未来证据与当前证据联合的候选框生成与筛选能力。在目标被遮挡、快速运动或相似物干扰导致当前帧证据缺失的情形下，方法能够依靠推演证据维持跟踪连续性；语言描述变更时，记忆与推演模块可随之重置，适配“改语言即换目标”的无模板特性。

性能目标：在TNL2K、LaSOT、OTB99\_Lang与MGIT四个评测基准上，取得较基线MVLM可测量的提升，其中遮挡与快速运动属性子集上的增益应显著高于整平均水平；若受基线复现条件限制而无法达成全部指标，则以“同一参照系下的相对提升”作为主要评价方式，确保增量结论可信。

成果目标：完成毕业论文一篇与可复现的代码实现，形成“内容记忆库、潜空间推演器、权值引导筛选”的完整方法方案。

\section{课题已具备和所需的条件、经费}

本课题已具备文献、代码、数据与算力条件，仅需少量经费支持备用算力租用与材料打印，具体情况如表~\ref{tab:t6-1}~所示。

\begin{table}[htbp]
  \caption{课题已具备条件与所需经费一览}
  \label{tab:t6-1}
  \centering
  \begin{tabularx}{\textwidth}[t]{A{0.20} A{0.60} B{0.20}}
    \toprule
    项目 & 内容 & 状态与需求 \\
    \midrule
    文献资料 & 课题方向相关文献与综述，含37篇参考文献（外文36篇） & 已具备 \\
    代码与工具 & 基线MVLM开源实现与复现、评测与统计脚本 & 已具备 \\
    数据资源 & TNL2K、LaSOT、VastTrack、OTB99\_Lang、MGIT等公开基准 & 已具备 \\
    算力条件 & 现有GPU资源可满足训练与评测需求 & 已具备 \\
    经费需求 & 备用云端算力租用、论文材料打印 & 少量 \\
    \bottomrule
  \end{tabularx}
  \tabdepthfix
\end{table}

\section{研究过程中可能遇到的困难和问题，解决的措施}
\label{sec:difficulty}

结合第~\ref{sec:content}~章提炼的三个关键问题与第~\ref{sec:route}~节的三阶段技术路线，预计研究过程中主要存在以下三方面困难与问题，本章逐项分析其成因与影响，并给出相应的解决措施。

\subsection{基线复现对齐困难}

MVLM\cite{park2026mvlm}的官方代码、配置与预训练权重虽已公开，复现基础较好，但复现指标与原文的差距仍存在超出第~\ref{sec:repro}~节所设一个百分点阈值的风险。造成差距的可能来源包括：训练数据预处理与数据增强细节未完整披露、预训练权重版本不一致、分布式训练与学习率调度策略难以精确还原，以及评测脚本在统计口径上的细微差别。若复现不对齐，后续记忆库与世界模型的增量收益将失去可信的参照系，全部对照实验的有效性都会受到质疑。

对此，将按以下步骤控制复现风险：首先，严格按官方配置与原文超参数搭建环境，先用官方权重跑通评测链路并核对指标，必要时通过邮件或开源社区与作者确认关键实现细节；其次，固定多组随机种子重复训练与评测，以均值±标准差报告结果，排除随机波动的影响；再次，若指标仍无法对齐，则按视觉编码器、跨模态融合层、预测头的顺序逐模块替换为官方实现，定位差异来源；最后，若个别差异确属原文未披露且无法还原，则以自复现基线作为全部对照实验的统一口径，并在论文中如实报告复现设置，保证增量对比在同一参照系下公平进行。

\subsection{记忆污染与记忆更新失衡}

内容记忆库的有效性依赖“写入的是有效特征”，但在遮挡、相似物干扰与误检等情形下，跟踪器可能将错误的外观与语义表征写入记忆库；错误条目会在后续帧中被注意力机制反复检索、反复强化，形成“污染—强化—再污染”的正反馈回路，误差随序列长度不断累积，最终导致跟踪漂移甚至目标丢失。同时，记忆更新策略面临两难：淘汰过激进会丢失仍有效的历史信息，削弱记忆对长时序跟踪的支撑；淘汰过保守则让过期与错误条目持续占用有限容量，引入噪声。此外，当语言描述变更、跟踪目标切换时，旧目标残留的记忆条目可能串扰到新目标的检索结果中。

对此，从写入、更新、重置三个环节协同治理：写入环节，以跟踪置信度与推演一致性双重门控筛选候选条目，只有当前帧证据可靠时才允许入库，从源头抑制错误特征；更新环节，采用冗余度淘汰与时间衰减相结合的自适应机制，被频繁检索命中的条目得以保留，长期未命中或与其他条目高度冗余的条目优先淘汰，兼顾信息留存与容量效率；重置环节，沿用第~\ref{sec:mem}~节的设计，语言描述变更时整体重置记忆，杜绝旧目标信息跨目标串扰；并在表~\ref{tab:t4-1}~记忆库组（M2）的消融框架下，以遮挡与长序列属性段上的实验数据驱动淘汰策略的最终选定。

\subsection{世界模型幻觉与推演偏差}

潜空间推演是本课题最核心也最具风险的模块：世界模型可能生成与目标真实状态不符的“幻觉”表征，一次错误推演即会把候选框引向错误位置，直接造成跟踪丢失。幻觉的来源主要有三：其一，推演条件中的记忆读出或语言语义本身存在偏差，错误条件被放大为错误预测；其二，仅有目标框标注、缺乏未来状态的直接监督，推演器只能依赖基于现有时序标注构造的回归目标间接学习，训练难度大；其三，若采用多步自回归展开，单步推演偏差会随步数累积，长时序推演逐渐偏离真实状态分布。

对此，从训练目标、推理机制与运行监控三个层面抑制幻觉：训练上，视觉前瞻仅以未来帧真值框内特征为回归目标、权值前瞻仅以下一帧候选框的实际重叠度与得分为回归目标，不引入易漂移的长程自回归展开，从机制上控制误差累积；推理上，以推演结果与实际观测的余弦相似度为门控，一致性高时前瞻证据正常参与候选框生成与筛选，一致性低时自动降权或弃用，保证推演“只在可靠时生效”；监控上，在验证集上持续跟踪幻觉度量指标，将其纳入模型选择与早期预警，并结合第~\ref{sec:ablation}~节消融实验量化前瞻证据在不同属性段上的实际贡献，确保推演模块带来的是可验证的正向收益。

\section{主要参考文献}

\bibliographystyle{hithesis}
\bibliography{reference}

% Local Variables:
% TeX-master: "../report"
% TeX-engine: xetex
% End:
